% HD-I17. GPT-6 Astra Pro, inherited effort, /root/imprimitive.
% 21 November 2026. Fresh interval verification, with full analytic proof.
\subsection{HD-I17: exact Fuchsian existence, conjugacy and the common residue box}
\label{sec:HD-I17}

This section verifies the auxiliary rank-eight Fuchsian system used
in HD-G17.  The numerical centre and target integer matrices come
from Benjamin Severn's candidate archive
\texttt{qi-m3-certificates}, supplied with an unrefereed manuscript.
They are used as explicit data.  No existence or rank theorem in
that manuscript is an input here.  The following fresh computation
and analytic argument prove the required existence and identify
its exact monodromy and its residue enclosure.

\paragraph{Statement and specified data.}
Let $a,b,c,d$ be the integer matrices in \texttt{mtuple.json},
datum \texttt{g}, tuple index $0$.  There is a Fuchsian system
\begin{equation}
 \frac{dY}{dz}=
 \left(\frac{B_0}{z}+\frac{u_1v_1^t}{z-1}
                         +\frac{u_2v_2^t}{z-(2+i)}\right)Y,
 \qquad
 B_0=\frac1{20}\operatorname{diag}(1,3,7,9,11,13,17,19),
                                                        \tag{I17.1}
\end{equation}
with $v_1^tu_1=v_2^tu_2=1/2$, and an invertible matrix $Q$,
whose based monodromies at $0,1,2+i$ are respectively
$QaQ^{-1},QbQ^{-1},QcQ^{-1}$.  For the explicitly fixed loops
below, the positive infinity monodromy is $QdQ^{-1}$.
Its residue
$B_\infty=-(B_0+u_1v_1^t+u_2v_2^t)$ is semisimple with exact
eigenvalue multiset
\begin{equation}
 -1,-1,-\frac{11}{12},-\frac23,-\frac7{12},
               -\frac5{12},-\frac13,-\frac1{12}.       \tag{I17.2}
\end{equation}
The unknown parameters lie in the exact archived residue box
used by the separate Laurent calculation HD-G17.

Here are the literal input identifiers.  The archive's centre
pickle contains only a list of $87$ pairs of decimal strings;
its gauge pickle has only lists, integers and the global
\texttt{builtins.complex}.  All opcodes were inspected before
conversion, and a restricted unpickler permitted that one
global only.  The converted JSON preserves every binary gauge
constant.  The SHA-256 values are:
\begin{center}\begin{minipage}{0.96\linewidth}\scriptsize\ttfamily
xtilde\_g\_0.pkl\\
afc056d5bb122e3c6d92c8dfbf6d5c54e2e8bd64468f875c2509ff08d8ce0649\\
meta\_g\_0.pkl\\
d60a5a641dbaef02a9bc02a0f4c4d9d16146f1ed032fe5d0408d884bebbbdb67\\
mtuple.json (archive and local copy)\\
239f3df1d88ba10ad4e805180040250447ba4e1a5b40a6cbc774c3a8c44be49d\\
center\_g.json (data-only conversion)\\
7c568f5df6c10702c1599ea60e47a8e35f0167c47c87decb15c5062392a263ae
\end{minipage}\end{center}
The complete script and transport-library hashes are in
\texttt{geometry\_and\_link.json}.  That file also records the
coordinatewise enclosure comparisons and all infinity minors.

\paragraph{1. Rigorous transport and the holomorphic equations.}
Fix $z_0=1/2-i$.  For each finite pole $p$, the path goes straight
from $z_0$ to the first vertex of a positively oriented regular
$12$-gon of radius $0.45$ centred at $p$, traverses its vertices,
and returns to $z_0$.  The vertex coordinates are the specific
binary rationals produced by the script; subsequent certification
treats these as exact rational points.  Their exact geometric
checks are described in section~4.

The eight coordinates of $u_1$ and one coordinate $u_{2,j_0}$
are fixed dyadic constants from the gauge data.  The variables are
the eight entries of $v_1$, the other seven entries of $u_2$,
the eight entries of $v_2$, and the $64$ entries of $Q$, for a
total of $87$.  Let $M_0,M_1,M_c$ denote transport around the
three loops.  These are holomorphic functions of all variables:
Picard's iterated-integral series along a fixed compact path
converges uniformly on compact parameter sets, bounded by the
exponential series for the supremum norm of its coefficient
matrix.  Thus it defines a holomorphic matrix function.

The transport enclosure used here has an explicit majorant.
For $A(z)=\sum A_k/(z-p_k)$ at a regular centre $z_*$, put
$\delta=\min_k|z_*-p_k|$ and
$S=\sum_k\|A_k\|_\infty$, using outward bounds on both.
The Taylor coefficient of $A$ in degree $j$ has norm at most
$S/\delta^{j+1}$.  The solution with value $I$ at the centre
is therefore coefficientwise majorized by
$(1-t/\delta)^{-S}$.  At a step $h$, set $q=|h|/\delta<1/2$.
After truncation at degree $N_0$, the tail is bounded by
\[
 \sum_{j>N_0}\binom{S+j-1}{j}q^j
 \leq
 \frac{\binom{S+N_0}{N_0+1}q^{N_0+1}}
 {1-q(1+S/(N_0+2))}.
\]
The denominator is checked positive.  The recurrence for the
matrix coefficients and this scalar bound are evaluated by
Arb complex balls; the tail is added to every matrix entry.
For boxes in the parameters the stable direct Taylor recurrence
is used, so all parameter dependence remains enclosed.
All steps end at their prescribed exact endpoints; the repaired
library asserts exact equality before skipping a purported
zero-length segment.  The library was read in full independently.
Arb's arithmetic semantics are its inclusion guarantee; see the
official \href{https://python-flint.readthedocs.io/en/latest/acb.html}
{\texttt{acb} documentation} and the
\href{https://flintlib.org/doc/arb.html}{FLINT real-ball documentation}.

Exact integer matrix calculations give
\[
 abcd=I,\quad a^{20}=I,\quad b^2=c^2=I,\quad
 \operatorname{rank}(I-b)=\operatorname{rank}(I-c)=1,
 \quad d^{12}=I,
\]
and
\[
 \det(xI-a)=x^8-x^6+x^4-x^2+1,\qquad
 \det(xI-d)=(x-1)^2(x^2+x+1)(x^4-x^2+1).
\]
Let $e_j$ be eigenvectors of $a$ for
$\lambda_j=\exp(2\pi i\beta_j)$, with $\beta_j$ the diagonal
entries of $B_0$.  Seven-row pivot solves with certified nonzero
minors compute enclosures for these exact vectors.  The exact
characteristic polynomial ensures that the omitted equation
holds.  Write $b=I-x_by_b^t$ and $c=I-x_cy_c^t$ by exact
rational factorizations, with $y_b^tx_b=y_c^tx_c=2$; choose exact
bases of the two seven-dimensional kernels of $y_b^t,y_c^t$.

Define the holomorphic map $F:\mathbf C^{87}\to\mathbf C^{87}$
by the following $87$ scalar equations: the two equations
$v_i^tu_i-1/2$; seven selected rows of each of the eight vectors
$(M_0-\lambda_j I)Qe_j$; seven selected rows of
$(M_1+I)Qx_b$ and seven single-pivot-row equations for
$(M_1-I)Qw$, as $w$ runs through the kernel basis; the analogous
$14$ equations for $M_c,c$; and $\det Q-1$.
The row choices are fixed once from the centre.  Every rank
assertion needed to justify these choices is subsequently
checked by ball minors on the whole solution enclosure.

\paragraph{2. The contraction proof and the fresh bounds.}
Let $\widetilde x$ be the $256$-bit midpoint of the data-only
decimal centre.  The new script
\texttt{verify\_fuchsian\_existence.py} evaluates $F$ there and
constructs an approximate derivative matrix $D$ as follows.
For the $23$ residue variables it uses the exact difference
quotient with $h=2^{-80}$, enclosing both function values.
For the other $64$ variables it computes the exact derivative
in the entry $Q_{uv}$: a term $MQe$ contributes $M_{iu}e_v$,
and the determinant contributes $\operatorname{adj}(Q)_{vu}$.
Thus these columns require no finite-difference estimate.
An ordinary floating inverse proposes a dyadic matrix $Y$;
all assertions about it are then checked by ball arithmetic.

Let $R_0=2^{-40}$ and let $K_i$ bound $|F_i|$ on the complex
polydisc of radius $R_0$ about $\widetilde x$.  The evaluated
rectangular input box contains that polydisc.  Cauchy's estimate
for a one-variable difference quotient gives
\[
 |D_{ij}-\partial_jF_i(\widetilde x)|
 \leq \frac{K_i h}{R_0(R_0-h)}
\]
in the first $23$ columns.  Hence the row-norm error is bounded by
$E=\max_i23K_i h/[R_0(R_0-h)]$.
On a polydisc of radius $\rho<R_0$, Cauchy's second-derivative
estimate gives the derivative variation bound
\[
 C=\max_i\frac{2\cdot87^2\rho K_i}{(R_0-\rho)^2}.
\]
The factor $2$ covers a repeated partial derivative; summing
over the moving variable and then the derivative's row gives
the two factors $87$.

Put $\eta=\|YF(\widetilde x)\|_\infty$ and choose the outward
dyadic value $\rho=10\eta$.  The fresh computation proved
\begin{align*}
 \max_iK_i&<1.607\cdot10^{-6},&
 \|I-YD\|_\infty&<7.14\cdot10^{-10},\\
 \|Y\|_\infty&<55.56,&
 \eta&<1.284\cdot10^{-35},\\
 \kappa:=\|I-YD\|_\infty+\|Y\|_\infty(E+C)
        &<0.002053,&
 \eta+\kappa\rho&<\rho.
\end{align*}
The certificate file \texttt{existence\_\allowbreak certificate.json}
records the full outward enclosures.
They were recomputed in $185$ seconds from the centre and gauge
data; the missing numerical starting solution and the missing
\texttt{fuchs.py} from the original runner are not used.

For $G(x)=x-YF(x)$ these inequalities give
$\|dG\|_\infty\leq\kappa<1$ on the closed polydisc of radius
$\rho$ and $\|G(x)-\widetilde x\|_\infty\leq\eta+\kappa\rho<\rho$.
The contraction principle gives a unique fixed point $x_*$ there.
Moreover $Y$ is invertible: the strict bound $\|I-YD\|<1$ makes
$YD$ invertible for the actual difference/derivative matrix
enclosed by $D$, and these are square matrices.  Thus the fixed
point satisfies $F(x_*)=0$, proving exact existence.
Its distance from the centre satisfies the stronger bound
\begin{equation}
 \|x_*-\widetilde x\|_\infty
 \leq\frac{\eta}{1-\kappa}
 <1.285978034319198\cdot10^{-35}.                 \tag{I17.3}
\end{equation}

\paragraph{3. Every omitted equation and exact conjugacy.}
At the root, $B_1=u_1v_1^t$ and $B_c=u_2v_2^t$ have eigenvalues
$1/2,0,\ldots,0$ and are semisimple.  The residue $B_0$ is diagonal
with the specified distinct eigenvalues.  No two eigenvalues of
any of these residues differ by a nonzero integer.  The local
holomorphic-gauge argument in section~5 therefore identifies
local monodromy with the exponential of the residue.
Consequently $M_0-\lambda_jI$ has rank seven,
$M_1+I,M_c+I$ have rank seven, and
$M_1-I,M_c-I$ have rank one.

On the whole parameter enclosure of radius $\rho$, the fresh
certificate verifies a nonzero $7\times7$ minor in each selected
seven-row block of $M_0-\lambda_jI$ and of $M_1+I,M_c+I$.
The seven equations defining the negative-eigenvector images
therefore imply their omitted eighth equation at the root:
the selected rows span the full rank-seven row space.
For each reflection, the chosen row of $M-I$ has a certified
nonzero entry.  At the root its rank is one, so vanishing of
that row on $Qw$ implies $(M-I)Qw=0$.
The vectors $x_b$ and the kernel basis of $y_b^t$ together form
a basis, because $y_b^tx_b=2$; the same holds for $c$.
Thus the equations prove
$M_0Q=Qa$, $M_1Q=Qb$, and $M_cQ=Qc$ on full bases.
Finally $\det Q=1$ makes these actual conjugacies.

\paragraph{4. Based-loop convention and the explicit pullback tuple.}
The exact geometry script uses rational cross products of the
binary vertices.  It checks that each closed $12$-gon is convex,
positively oriented and contains its own pole, while all its
vertices lie within distance $0.46$ of that pole.  The other poles
are outside that disk.  It also checks that the small triangles
used to replace the approximate final vertex by the first and
to straighten the stem contain no puncture.  Hence each loop
is the positive meridian with straight radial stem from $z_0$.

The stem directions, in counterclockwise order, are $c,1,0$.
The three exact cross products for this order are $1/2,5/2,1$,
all positive, and all stems lie in the upper half-plane from
$z_0$.  Following the positively oriented boundary of a small
regular neighbourhood of this embedded tree traverses the
meridians in the path order $\gamma_c,\gamma_1,\gamma_0$.
For $Y'=A Y$, transport after the first path is multiplied on
the left by transport along the next path.  Therefore the
positive boundary about all finite poles has matrix
$M_0M_1M_c=QabcQ^{-1}$, and the positive infinity meridian
has its inverse $QdQ^{-1}$.
This order matters: direct exact computation shows that the
reversed matrix product $cba$ has a different characteristic
polynomial, so it is not substituted for $abc$.

For later use, the actual local system
$s^{1/4-3}M(s^5)$ can be described by lifted meridians with
the following product-one tuple:
\begin{equation}
 \left(i a^5,
  a^{-4}ba^4,a^{-4}ca^4,
  a^{-3}ba^3,a^{-3}ca^3,\ldots,b,c,-i d^5\right).
                                                        \tag{I17.4}
\end{equation}
These are the successive lifts represented by the five
Schreier sheets about $0$, with the scalar character $i$
at $0$ and $-i$ at infinity.  Their product is checked by the
identity
\[
 a^5\prod_{j=4,3,2,1,0}a^{-j}(bc)a^j=(abc)^5,
\]
followed by $(abc)^5d^5=I$.  HD-G17 uses this precise tuple
and computes its group and cup form directly; no identification
with a reversed or transposed archived tuple is assumed.

\paragraph{5. The exact infinity exponents.}
An approximate eigenbasis of the enclosed $B_\infty$ is used
only to propose a dyadic invertible matrix $V$.  Its determinant
is certified nonzero.  For $T=V^{-1}B_\infty V$, the new verifier
uses the outward Arb bound
\[
 \left|T_{jj}-q_j\right|+\sum_{k\ne j}|T_{jk}|
        <2.874\cdot10^{-15},
\]
where the proposed rational centres $q_j$ have exactly the
multiset (I17.2).  This replaces the original runner's
floating-point radius printout.  The Gershgorin components
are pairwise disjoint apart from the two centred at $-1$;
they contain the indicated numbers of eigenvalues.

All eigenvalues therefore lie in a strip of width less than one,
and no difference of two is a nonzero integer.  Here is the
local gauge fact used both here and in section~3.  For a
regular-singular equation with residue $B$, seek a solution
$H(t)t^B$, $H(0)=I$.  If the analytic remainder is
$\sum_{j\geq0}A_jt^j$, the coefficient recurrence is
\[
 (kI-\operatorname{ad}B)H_k
          =\sum_{j=0}^{k-1}A_jH_{k-1-j}.
\]
The stated eigenvalue-difference condition makes all left
operators invertible.  Their inverse norms are at most $C/k$
for one finite $C$: for large $k$ this is the Neumann-series
bound, and the remaining finitely many $k$ enlarge $C$.
An analytic scalar majorant for
$h'=C A/(1-t/R)\,h$ proves convergence of $H$ on a smaller
disk.  Thus local monodromy is conjugate to $\exp(2\pi iB)$.

At infinity it is also conjugate to the exact matrix $d$,
which satisfies $d^{12}=I$.  The exponential is therefore
semisimple, and so is $B_\infty$: a nontrivial Jordan block
would retain a nontrivial nilpotent part under the exponential.
Each eigenvalue of $B_\infty$ is consequently a real multiple
of $1/12$ modulo integers.  Each certified disk contains
only its displayed such rational number.  This proves (I17.2)
exactly, not by rounding a numerical spectrum.

As a separate conditioning check, the geometry/link script
certifies a nonzero rank-seven minor of $B_\infty-qI$ for
each of the six simple eigenvalues, and a nonzero rank-six
minor for $q=-1$.  It records their actual rows, columns and
ball determinants; every determinant excludes zero.  This
connects the exact spectral multiplicities to the eigenvector
solves used by the Laurent calculation.

\paragraph{6. Literal linkage to the independent Laurent rank box.}
The archived radius used in HD-G17 is the exact binary value
represented by
$r_L=1.0528417230446196\cdot10^{-34}$.
The centre strings and fixed gauge constants are the same
inputs whose hashes appear above.  In addition to the strict
inequality in (I17.3), the new script checks all $87$ literal
rectangular inclusions
\[
 [\Re\widetilde x_j\pm\delta]
       +i[\Im\widetilde x_j\pm\delta]
 \ \subset\
 \operatorname{acb}\bigl(\operatorname{arb}(a_j,r_L),
                         \operatorname{arb}(b_j,r_L)\bigr),
 \qquad \delta=\operatorname{upper}\frac{\eta}{1-\kappa},
\]
where $(a_j,b_j)$ are the original decimal strings.  Every
inclusion passes, so the existing independent Laurent
calculation applies to the exact Fuchsian system proved here.

\paragraph{Audit and reproducibility.}
The full source \texttt{cert\_M.py} and transport library were
read before constructing the fresh verifier.  The new verification
uses the original explicit candidate centre, but independently
rebuilds its equations, pivot choices, derivative calculation,
contraction bounds and infinity enclosures.  It computes $64$
derivative columns exactly instead of repeating the original
finite differences.  The fresh interval run passed; the separate
exact geometry/box/minor run passed.  The author, GPT-6 Astra Pro
with inherited effort, read the complete analytic argument and
its relation to both computations on 21 November 2026.  These are
distinct events; a successful numerical run is not described
as a substitute for the omitted-row, gauge or topological proofs.
The files are in \texttt{work/imprimitive/e8/}; the two result
files are \texttt{existence\_certificate.json} and
\texttt{geometry\_and\_link.json}.
